\section{Future work}\label{sec:future-work}

% Mathematical job:
%   - List future research directions WITHOUT promising results not earned.
%   - Topics:
%       endogenous instrument meta-limit / self-certification theorem;
%       broader-domain exact-six stress beyond FATCD;
%       uniqueness or canonical decomposition study;
%       full six-primitives algebra;
%       empirical substrate mapping;
%       formal mechanization strengthening (Lean);
%       cross-instrument comparison of exact-six decompositions.
%
% Governing sources:
%   workspace/meta/lens_lab/papers/paper3/gaps_and_followups.md
%   vision.md (Stage 6: Empirical bridge; meta-limit notes)
%
% Overclaim risks:
%   - Writing future work as commitments or predictions.
%   - Reopening the terminal theorem by framing these as gaps.

Theorem~\ref{thm:scoped-exact-six} and the limitations of
Section~\ref{sec:limitations-nonclaims} suggest several follow-on
research directions. Each subsection below states a direction, not a
commitment.

\subsection{Endogenous-instrument meta-limit}\label{subsec:future-meta-limit}

A formalization of the meta-limit described in
Section~\ref{subsec:meta-limit} would target a three-way
impossibility statement: no internal argument can simultaneously
establish the scoped exact-six claim, certify the soundness of the
accepted instrument family, and certify its completeness for
exact-six judgments formulable inside the same regime.
Theorem~\ref{thm:scoped-exact-six} would remain intact under such a
result; the additional content would be a precise account of
self-certification boundaries.

\subsection{Broader-domain exact-six}\label{subsec:future-broader-domain}

Candidate extensions beyond \FATCD{} include admissibility regimes
that weaken finiteness, closure, or audit
(Section~\ref{sec:fatcd-domain}), and regimes that admit annotation
kinds beyond the six \(\mathrm{Ann}_*\) categories of
Definition~\ref{def:annotations}. Any such extension would require
fresh upper-bound and decomposition/exhaustion analysis, since the
alignment rules and threat closures are not stated at scope-free
strength.

\subsection{Uniqueness or canonical decomposition}\label{subsec:future-uniqueness}

Theorem~\ref{thm:scoped-exact-six} is existential. A uniqueness study
would investigate when distinct well-formed decomposition/exhaustion
records of the same \(D \in \FATCD{}\) are equivalent under explicit
translation records, or when a canonical representative can be
selected. The channel-versus-activation discipline
(Proposition~\ref{prop:channel-vs-activation}) is compatible with
such a study; no uniqueness claim is part of the present theorem.

\subsection{Full six-primitives algebra}\label{subsec:future-algebra}

The generator-and-relation structure of \Pone{}--\Psix{} beyond
typed non-collapse remains open. A full algebra would require
scope-free independence data, not just the local boundary
distinctions of Table~\ref{tab:boundary-matrix}, and would likely
require a categorical apparatus richer than the raw grammar of
Section~\ref{subsec:raw-grammar}.

\subsection{Empirical substrate mapping}\label{subsec:future-substrate}

The empirical-bridge admissibility regime
(Proposition~\ref{prop:empirical-bridge}) can be connected to
concrete substrate-level observables. The present paper treats
empirical bridges as threshold-dependent, level-indexed, and
instrument-visible annotations; later work can couple this to
substrate-specific measurement theory while preserving the nonclaim
of empirical realization.

\subsection{Mechanization}\label{subsec:future-mechanization}

The typed non-collapse proposition, the decomposition-existence
theorem, and the upper-bound classification theorem are natural
targets for mechanization in a proof
assistant~\citep{DeMouraUllrich2021}. Work-in-progress Lean
mechanization accompanying the present paper is available at the
code repository of Appendix~\ref{app:code-availability}, and the
Foundations program~\citep{Tsiokos2026SixBirdsFoundations} has
already mechanized closure-related infrastructure that can be
reused.

\subsection{Cross-instrument comparison}\label{subsec:future-cross-instrument}

Decomposition/exhaustion records of the same \(D \in \FATCD{}\) can
be compared under different admissible instrument assignments.
Proposition~\ref{prop:visibility} admits visibility records that
suppress different portions of a claim under different instruments; a
cross-instrument study could make the instrument-dependence of
\(\Dec{D}\) explicit and sharpen the notion of instrument-relative
claim strength.
